\chapter{Writing Mathematics from the Beginning}
\label{chap:math-beginning}
\index{mathematics}\index{equation}\index{math mode}\index{align environment}

Mathematics in a scientific document is part of a sentence. It must be
accurate, readable, and connected to the variables around it. \LaTeX{} gives
mathematics a separate mode because symbols such as a minus sign, a variable,
and a fraction need different spacing from ordinary prose.

\begin{learningoutcomes}
  \item Distinguish inline and displayed mathematics.
  \item Create scripts, fractions, roots, Greek letters, and functions.
  \item Size brackets around structured expressions.
  \item Align related equations with \pkg{amsmath}.
  \item Number, label, and refer to an equation.
  \item Avoid obsolete display syntax and manual spacing repairs.
\end{learningoutcomes}

\section{Inline mathematics belongs in a sentence}

Place short mathematics between \code{\$} delimiters:

\begin{latexcode}{Complete example: inline mathematics}
\documentclass{article}
\begin{document}
Let $d_i = B_i - A_i$ be the difference for pair $i$.
\end{document}
\end{latexcode}

The dollar signs switch mathematics on and off; they do not print. Letters in
math mode are treated as variables and normally appear italic. The minus sign
has mathematical spacing. Keep sentence punctuation outside or after the
formula as grammar requires.

\begin{commonerror}
\textbf{Symptom:} the rest of a paragraph appears in mathematical italics or
the compiler reports a missing dollar sign.

\textbf{Cause:} an opening \code{\$} has no matching closing delimiter.

\textbf{Repair:} inspect the equation and the source just before the reported
line. Match the delimiters, then compile again.
\end{commonerror}

\section{Displayed mathematics}

Use \code{\textbackslash[} and \code{\textbackslash]} for an unnumbered
display:

\begin{latexcode}{Unnumbered display}
The arithmetic mean is
\[
  \bar{x} = \frac{1}{n}\sum_{i=1}^{n} x_i.
\]
\end{latexcode}

The display is still part of the sentence, so it ends with a full stop.
Indentation is for source readability.

Do not use \code{\$\$ ... \$\$}. That plain-TeX form bypasses \LaTeX{}'s display
management and can give inconsistent spacing. Numbered displays use an
appropriate environment, not a manually typed number.

\section{Scripts need groups}

The caret creates a superscript and the underscore a subscript in mathematics:

\begin{latexcode}{Scripts and grouping}
$x^2$, $x_i$, $x_{ij}$, $10^{-3}$, and $a^{n+1}$
\end{latexcode}

Without braces, only the next token belongs to the script. Thus
\code{x\_ij} makes only \code{i} a subscript and leaves \code{j} at the main
level. Group multi-character scripts: \code{x\_\{ij\}}.

\begin{tryit}
Compile \code{\$x\_ij\$} and \code{\$x\_\{ij\}\$} side by side. Describe the
baseline of the letter \code{j} in each output.
\end{tryit}

\section{Fractions, roots, and Greek letters}

\cmd{frac} takes a numerator and denominator. \cmd{sqrt} takes the expression
under a root and accepts an optional root index.

\begin{latexcode}{Complete example: common constructions}
\documentclass{article}
\usepackage{amsmath}
\begin{document}
\[
  z = \frac{x - \mu}{\sigma}, \qquad
  r = \sqrt{x^2 + y^2}, \qquad
  q = \sqrt[3]{V}.
\]
\end{document}
\end{latexcode}

Commands such as \cmd{mu}, \cmd{sigma}, \cmd{alpha}, and \cmd{Delta} insert
Greek letters. Upper- and lower-case command names can represent different
symbols. Do not paste a look-alike Unicode character into an otherwise ASCII
formula without checking the engine and font support.

The \code{\textbackslash qquad} spacing in this compact demonstration separates
three independent expressions. It should not be used to repair the internal
structure of one badly expressed formula.

\section{Functions and operators}

Standard functions have upright names and mathematical spacing:

\begin{latexcode}{Functions and operators}
$\sin \theta$, $\log x$, $\exp(-kt)$, and $\max_i x_i$
\end{latexcode}

Use the supplied commands rather than typing italic letters such as
\code{sin}. For a recurring named operator not supplied by \LaTeX{},
\pkg{amsmath} provides \cmd{DeclareMathOperator}; define it once in the
preamble after understanding its mathematical role.

Ordinary words inside mathematics use \cmd{text} from \pkg{amsmath}:

\begin{latexcode}{Words inside a formula}
\[
  y = \begin{cases}
    1, & \text{if the pair is complete},\\
    0, & \text{otherwise}.
  \end{cases}
\]
\end{latexcode}

The ampersand marks the alignment point and \code{\\} ends a row inside this
structured environment. They are not paragraph-formatting tricks.

\section{Brackets that fit}

Ordinary parentheses do not grow automatically around a tall fraction.
\cmd{left} and \cmd{right} request paired delimiters sized to their contents:

\begin{latexcode}{Sized delimiters}
\[
  \left(\frac{a+b}{c+d}\right)^2
\]
\end{latexcode}

Every \cmd{left} needs a matching \cmd{right}. When only one visible delimiter
is required, a full stop can act as an invisible partner, but that technique
should be introduced only with a real notation need. The \pkg{mathtools}
package also supports reusable paired-delimiter commands for larger projects.

\section{Numbering and referencing an equation}

The \code{equation} environment numbers one display. Place the label after the
equation begins and give it a descriptive prefix:

\begin{equation}
  \bar{d}=\frac{1}{n}\sum_{i=1}^{n}(B_i-A_i).
  \label{eq:mean-difference}
\end{equation}

The source is:

\begin{latexcode}{Numbered and labelled equation}
\begin{equation}
  \bar{d}=\frac{1}{n}\sum_{i=1}^{n}(B_i-A_i).
  \label{eq:mean-difference-example}
\end{equation}

Equation~\eqref{eq:mean-difference-example} defines the mean difference.
\end{latexcode}

\cmd{eqref} prints the equation number with parentheses. Compilation may need
a repeat pass after a new label. A label is a stable name; the visible number
can change when an earlier equation is inserted.

\begin{scienceinpractice}
In the running comparison, $A_i$ and $B_i$ are paired simulated values and
$\bar{d}$ estimates their mean signed difference. The equation does not by
itself establish agreement; the surrounding methods must define sampling,
units, missingness, and interpretation.
\end{scienceinpractice}

\section{Aligned equations}

Use \code{align} for related lines. The ampersand marks the alignment point,
usually before a relation such as equals:

\begin{latexcode}{Complete example: aligned equations}
\documentclass{article}
\usepackage{amsmath}
\begin{document}
\begin{align}
  d_i &= B_i-A_i, \\
  \bar{d} &= \frac{1}{n}\sum_{i=1}^{n}d_i.
\end{align}
\end{document}
\end{latexcode}

Use \code{align*} when none of the lines should be numbered. Use
\cmd{notag} selectively when a single line in a numbered alignment should not
receive a number.

Do not use the older \code{eqnarray} environment. Its spacing around relation
symbols is inferior and its control is limited. Do not align equations by
inserting guessed spaces; state the alignment point structurally.

\begin{goodpractice}
Define every symbol near its first use, keep notation consistent, and inspect
the meaning as well as the typesetting. A perfectly aligned equation with an
undefined variable is not publication-ready.
\end{goodpractice}

\chapterproject

\begin{miniproject}
Load \pkg{amsmath} and \pkg{mathtools} in the running article. In Methods,
define $d_i=B_i-A_i$. Add a labelled equation for the mean difference and refer
to it from the prose with \cmd{eqref}. State clearly that the values are
simulated and that the sign depends on the chosen subtraction order.
\end{miniproject}

\chaptersummary

Dollar delimiters create inline mathematics; \code{\textbackslash[} and
\code{\textbackslash]} create an unnumbered display. Braces group scripts and
arguments. Mathematical commands produce fractions, roots, symbols, functions,
text, and sized delimiters. \code{equation} numbers one display; \code{align}
organises related lines. Labels and \cmd{eqref} replace typed numbers.
Avoid \code{\$\$}, \code{eqnarray}, and manual spacing repairs.

\chaptercheck
\begin{chapterchecklist}
  \item Every math delimiter and brace has a partner.
  \item Multi-character scripts are grouped.
  \item Function names are upright commands.
  \item Display punctuation follows the surrounding sentence.
  \item Equation references come from labels.
  \item Every symbol in the running article is defined.
\end{chapterchecklist}

\chapterexercisestitle
\begin{chapterexercises}
  \item Write inline source for $E=mc^2$.
  \item Correct \code{x\_ij} so both letters form the subscript.
  \item Typeset the fraction $(x-\mu)/\sigma$ with a built-up fraction.
  \item Replace italic \code{sin} with the mathematical function command.
  \item Create two equations aligned at their equals signs.
  \item Label an equation and refer to it without typing its number.
  \item Explain why \code{\$\$} and \code{eqnarray} are discouraged.
  \item Challenge: typeset and define one equation from your discipline, then
        ask a colleague whether the notation is unambiguous.
\end{chapterexercises}
